\documentclass[../general/general]{subfiles}

\begin{document}

\chapter{Лекция №5}
\noindent\textit{А.\,М.~Белемук \hfill 13 марта 2021}
\vspace{1em}


% ===== СЕКЦИЯ 1: Термодинамические следствия Большого канонического ансамбля =====
\section{Термодинамические следствия Большого канонического ансамбля}

Так, сегодня у нас лекция пятая, статистическая физика. В прошлый раз мы с вами ввели большой канонический ансамбль, большое каноническое распределение. И мы с вами показали, что был введен Омега-потенциал ($\Omega$), и показали, что это есть функция от температуры $T$, объема $V$ и химического потенциала $\mu$:
\begin{equation}
    \Omega = \Omega(T, V, \mu)
\end{equation}

Кроме того, мы еще показали, что Омега-потенциал представляется в виде:
\begin{equation}
    \Omega = -P(T, \mu) V
\end{equation}
где $P$ — давление. Давайте сразу отсюда выведем полезные соотношения. 

\subsection{Уравнение Гиббса-Дюгема из $\Omega$-потенциала}

Во-первых, мы знаем, что дифференциал Омега-потенциала $d\Omega$ равен:
\begin{equation}
    d\Omega = -S dT - P dV - N d\mu
\end{equation}
А теперь давайте сравним это с дифференциалом от правой части, то есть возьмем дифференциал от выражения $-PV$:
\begin{equation}
    d(-PV) = -P dV - V dP
\end{equation}
И это все равно $d\Omega$. 

Сравниваем два равенства между собой. Они должны быть тождественно равны. Мы видим, что член $-P dV$ присутствует с обеих сторон, поэтому он сокращается, и остается:
\begin{equation}
    -V dP = -S dT - N d\mu \implies V dP = S dT + N d\mu
\end{equation}

Следовательно, я вижу, что дифференциал давления можно выразить следующим образом:
\begin{equation}
    dP = \frac{S}{V} dT + \frac{N}{V} d\mu
\end{equation}

\subsection{Плотности термодинамических величин}

Посмотрим внимательно на коэффициенты перед дифференциалами в полученном выражении. Отношение $S/V$ есть энтропия на единицу объема, то есть плотность энтропии, обозначим ее малой буквой $s$. Отношение $N/V$ есть плотность частиц (или концентрация частиц), обозначим ее $n$. Таким образом:
\begin{equation}
    dP = s dT + n d\mu
\end{equation}

Из этого равенства мы можем считать частные производные давления $P$, рассматриваемого как функция от $T$ и $\mu$. Частная производная давления по температуре при постоянном химическом потенциале дает плотность энтропии:
\begin{equation}
    \left( \frac{\partial P}{\partial T} \right)_{\mu} = \frac{S}{V} = s
\end{equation}
А частная производная давления по химическому потенциалу при постоянной температуре дает плотность частиц:
\begin{equation}
    \left( \frac{\partial P}{\partial \mu} \right)_{T} = \frac{N}{V} = n
\end{equation}
Мы получили плотность частиц как функцию температуры и химического потенциала. Видим, что дифференцированием давления по химическому потенциалу мы легко находим концентрацию. Давление как функция химического потенциала имеет вот такие полезные частные производные, которые зависят от $T$ и $\mu$. 

\subsection{Соотношения для термодинамического потенциала Гиббса ($\Phi$)}

Аналогично, давайте для полноты картины выведем еще одно полезное соотношение. Мы в прошлый раз на этом остановились. Термодинамический потенциал Гиббса, рассматриваемый как функция переменных $T, P$ и $N$, пропорционален экстенсивной величине $N$, а в качестве удельного термодинамического потенциала фигурирует химический потенциал $\mu$:
\begin{equation}
    \Phi(T, P, N) = \mu(T, P) N
\end{equation}

Снова берем и сравниваем дифференциалы. Мы помним, что дифференциал термодинамического потенциала Гиббса равен:
\begin{equation}
    d\Phi = -S dT + V dP + \mu dN
\end{equation}
И сравниваем с дифференциалом произведения $\mu N$:
\begin{equation}
    d(\mu N) = \mu dN + N d\mu
\end{equation}

Приравнивая одно другому, слагаемое $\mu dN$ у нас сокращается с обеих сторон, и мы получаем:
\begin{equation}
    N d\mu = -S dT + V dP
\end{equation}
Отсюда я нахожу дифференциал $d\mu$:
\begin{equation}
    d\mu = -\frac{S}{N} dT + \frac{V}{N} dP
\end{equation}

Здесь отношение $S/N$ — это удельная энтропия, приходящаяся на одну частицу. Отношение $V/N$ — это удельный объем, приходящийся на одну частицу (обозначим его $v$). Таким образом, химический потенциал $\mu$ как функция температуры и давления имеет следующие частные производные:
\begin{equation}
    \left( \frac{\partial \mu}{\partial T} \right)_{P} = -\frac{S}{N} = -s_\text{уд}
\end{equation}
\begin{equation}
    \left( \frac{\partial \mu}{\partial P} \right)_{T} = \frac{V}{N} = v
\end{equation}

Кроме того, из равенства $V dP = S dT + N d\mu$ мы видим, что переменные $\mu, T$ и $P$ не могут являться независимыми. Они зависимы друг от друга. Это соотношение в термодинамике называется соотношением Гиббса-Дюгема. То, что $\mu, T, P$ — зависимые переменные, это важный факт. 

Вот мы сделали необходимые дополнения к тому, на чем остановились в прошлый раз. Теперь давайте посчитаем флуктуации числа частиц в Большом каноническом ансамбле.

% ===== КОНЕЦ СЕКЦИИ 1 =====


% ===== СЕКЦИЯ 2: Флуктуации числа частиц в открытой системе =====
\section{Флуктуации числа частиц в открытой системе}

Каковы же флуктуации числа частиц в Большом каноническом ансамбле? В Большом каноническом ансамбле число частиц $N$ у нас не фиксировано, оно непрерывно флуктуирует около какого-то среднего значения $\langle N \rangle$. 

\subsection{Определение дисперсии числа частиц}

Я полагаю, что функция распределения по $N$ (вероятность найти систему с заданным числом частиц) будет иметь какой-то вид. Позже мы увидим, что эта функция будет иметь резкий пик около среднего значения $\langle N \rangle$. 

\subsection{Связь флуктуации со статистической суммой}

Вспомним полезную операцию, которую мы вводили в прошлый раз. Мы знаем, что оператор $T \frac{\partial}{\partial \mu}$, примененный к логарифму большой статсуммы $Z$, дает среднее число частиц $\langle N \rangle$:
\begin{equation}
    T \frac{\partial}{\partial \mu} \ln Z = \langle N \rangle
\end{equation}
Эта операция $T \frac{\partial}{\partial \mu}$ очень полезна. Давайте применим её к самому среднему числу частиц $\langle N \rangle$, чтобы посмотреть, что получится:
\begin{equation}
    T \frac{\partial \langle N \rangle}{\partial \mu} = T \frac{\partial}{\partial \mu} \left( \frac{1}{Z} \sum_{k,N} N e^{\frac{\mu N - E_{k,N}}{T}} \right)
\end{equation}

Я вычисляю эту производную. Дифференцирую первый множитель $1/Z$, получаю:
\begin{equation}
    - \frac{1}{Z^2} \frac{\partial Z}{\partial \mu} T \sum_{k,N} N e^{\frac{\mu N - E_{k,N}}{T}}
\end{equation}
Здесь мы замечаем, что $T \frac{1}{Z} \frac{\partial Z}{\partial \mu}$ — это $\langle N \rangle$, а оставшаяся часть $\frac{1}{Z} \sum N e^{\dots}$ — это тоже по определению $\langle N \rangle$. Значит, этот член дает просто $-\langle N \rangle \langle N \rangle = -\langle N \rangle^2$.

Теперь дифференцирую экспоненту. $1/Z$ остается, $N$ остается, производная от экспоненты по $\mu$ спускает множитель $N/T$. Температура $T$ перед производной и в знаменателе сокращается, остается множитель $N$. Итого получаем:
\begin{equation}
    + \frac{1}{Z} \sum_{k,N} N \cdot N \cdot e^{\frac{\mu N - E_{k,N}}{T}} = \frac{1}{Z} \sum_{k,N} N^2 e^{\frac{\mu N - E_{k,N}}{T}} = \langle N^2 \rangle
\end{equation}

\subsection{Итоговое выражение для дисперсии}

Собирая оба слагаемых вместе, мы получаем:
\begin{equation}
    T \frac{\partial \langle N \rangle}{\partial \mu} = \langle N^2 \rangle - \langle N \rangle^2
\end{equation}
Но разность среднего квадрата и квадрата среднего — это, по определению, средний квадрат флуктуации (дисперсия) $\langle \Delta N^2 \rangle$. Таким образом, мы вывели фундаментальную формулу:
\begin{equation}
    \langle \Delta N^2 \rangle = T \frac{\partial \langle N \rangle}{\partial \mu}
\end{equation}
Среднеквадратичная флуктуация числа частиц определяется производной от среднего числа частиц по химическому потенциалу.

\subsection{Относительная флуктуация и эквивалентность ансамблей}

Много это или мало? Чтобы это понять, нужно сравнивать абсолютную флуктуацию $\sqrt{\langle \Delta N^2 \rangle}$ со средним значением $\langle N \rangle$. Мы должны посмотреть на относительную флуктуацию:
\begin{equation}
    \frac{\sqrt{\langle \Delta N^2 \rangle}}{\langle N \rangle} \sim \frac{\sqrt{N}}{N} = \frac{1}{\sqrt{N}} \ll 1
\end{equation}
Поскольку производная $\frac{\partial \langle N \rangle}{\partial \mu}$ пропорциональна $N$ (число частиц экстенсивно, а $\mu$ интенсивно), числитель пропорционален $\sqrt{N}$. В пределе макроскопической системы ($N \to \infty$) относительная флуктуация стремится к нулю. 

Это означает, что функция распределения по числу частиц в системе $W(N)$ имеет исключительно узкий пик в окрестности $\langle N \rangle$.

\begin{figure}[h!]
    \centering
    \includegraphics[width=0.3\textwidth]{W(N)}
    \caption{Распределение вероятности $W(N)$ в Большом каноническом ансамбле.}
\end{figure}

Подавляющая часть систем в ансамбле имеет число частиц, практически не отличимое от среднего $\langle N \rangle$. В этом смысле, канонический ансамбль (то, что он предсказывает для средних величин) эквивалентен — ну, с оговорками, что вы не учитываете флуктуации числа частиц, так как они ничтожно малы — эквивалентен Большому каноническому ансамблю. Вы можете спокойно заменять точное число частиц на среднее и пользоваться удобным математическим аппаратом Большого канонического ансамбля.

Вот мы с теориями Большого канонического ансамбля очень хорошо все обсудили. Теперь тема нашей лекции переходит непосредственно к применению этого аппарата. Мы стартуем с идеального ферми-газа.

% ===== КОНЕЦ СЕКЦИИ 2 =====







% ===== СЕКЦИЯ 3: Идеальный квантовый газ: модель и Большой канонический ансамбль =====
\section{Идеальный квантовый газ: модель и Большой канонический ансамбль}

Тема нашей лекции переходит непосредственно к применению математического аппарата. Вот наша тема на сегодня: Большой канонический ансамбль для идеального ферми-газа.

\subsection{Модель невзаимодействующих частиц}

Давайте сначала посмотрим, что из себя представляет ферми-газ. Мы рассматриваем нашу систему, занимающую объем $V$, представим ее в виде ящика.


Внутри этого ящика есть одночастичные уровни энергии, или орбитали. Каждой орбитали соответствует своя волновая функция. Почему я говорю об одночастичных уровнях энергии? Потому что газ у нас идеальный, частицы не взаимодействуют друг с другом.

Что я могу написать для гамильтониана такой системы? Гамильтониан $\hat{H}$ — это просто сумма по всем $i$ от единицы до $N$ одночастичных гамильтонианов $\hat{h}_i$:
\begin{equation}
    \hat{H} = \sum_{i=1}^N \hat{h}_i
\end{equation}

Каждый одночастичный гамильтониан $\hat{h}_i$ равен оператору кинетической энергии $\frac{\hat{p}_i^2}{2m}$. Если есть какой-то внешний потенциал, то добавляется слагаемое $U(\vec{r}_i)$. Если присутствует внешнее магнитное поле, то есть еще зависимость от поля, например, слагаемое $-\hat{\vec{\mu}} \cdot \vec{H}$, описывающее взаимодействие магнитного момента электрона с внешним магнитным полем:
\begin{equation}
    \hat{h}_i = \frac{\hat{p}_i^2}{2m} + U(\vec{r}_i) - \hat{\vec{\mu}} \cdot \vec{H}
\end{equation}

Собственные функции гамильтониана $\hat{h}_i$ без внешнего потенциала (в пустом ящике) — это плоские волны $|\vec{p}\rangle$. Кроме того, надо учесть проекцию спина $\sigma$. То есть волновая функция одночастичного состояния задается вектором $|\vec{p}, \sigma\rangle$, который мы будем обозначать как $\psi_{\vec{p}, \sigma}$.

\subsection{Микросостояния через числа заполнения}

Так как газ невзаимодействующий, происходит простое заполнение этих одночастичных уровней энергии. Мы вводим числа заполнения $n_{p\sigma}$. Это число показывает, сколько частиц сидит на каждом одночастичном уровне энергии в выбранном нами базисе волновых функций. 

Для ферми-системы число заполнения $n_{p\sigma}$ принимает значение либо $0$, либо $1$:
\begin{equation}
    n_{p\sigma} = 0, 1
\end{equation}
Либо на данном состоянии ноль частиц, либо сидит ровно один фермион. 

Полное число частиц в системе $N$ полностью определяется этим набором. Набор всех чисел заполнения для всех состояний $\{n_{p\sigma}\}$ определяет конкретное микросостояние многочастичной системы. Полное число частиц равно сумме чисел заполнения по всем импульсам и спинам:
\begin{equation}
    N = \sum_{\vec{p}, \sigma} n_{p\sigma}
\end{equation}

А полная энергия данного микросостояния $E$ также зависит от набора $\{n_{p\sigma}\}$ и равна сумме произведений одночастичных энергий $\epsilon_{p\sigma}$ на соответствующие числа заполнения:
\begin{equation}
    E = \sum_{\vec{p}, \sigma} \epsilon_{p\sigma} n_{p\sigma}
\end{equation}
Это справедливо для любого конкретного чистого микросостояния невзаимодействующей системы.

\subsection{Большая статистическая сумма идеального газа}

Далее нам нужно найти Большую статистическую сумму для нашего ферми-газа. По определению, Большая статистическая сумма $\mathcal{Z}$ задается двойной суммой:
\begin{equation}
    \mathcal{Z} = \sum_{N=0}^{\infty} \sum_{k} \exp\left( \frac{\mu N - E_{k,N}}{T} \right)
\end{equation}
Здесь индекс $k$ нумерует квантовые состояния при заданном числе частиц $N$. Но индекс $k$ — это, по сути, просто синоним конкретного набора чисел заполнения $\{n_{p\sigma}\}$. Энергия $E_{k,N}$ и число частиц $N$ полностью определяются этим набором.

Поэтому выражение для статсуммы можно переписать иначе: это просто сумма по всем возможным наборам чисел заполнения. Она автоматически включает в себя как перебор всех квантовых состояний при заданном $N$, так и перебор всех возможных чисел $N$.
\begin{equation}
    \mathcal{Z} = \sum_{\{n_{p\sigma}\}} \exp\left( \frac{\mu \sum_{\vec{p}, \sigma} n_{p\sigma} - \sum_{\vec{p}, \sigma} \epsilon_{p\sigma} n_{p\sigma}}{T} \right)
\end{equation}
Мы можем объединить суммы в показателе экспоненты, вынеся $n_{p\sigma}$ за скобки:
\begin{equation}
    \mathcal{Z} = \sum_{\{n_{p\sigma}\}} \exp\left( \sum_{\vec{p}, \sigma} \frac{(\mu - \epsilon_{p\sigma}) n_{p\sigma}}{T} \right)
\end{equation}

% ===== КОНЕЦ СЕКЦИИ 3 =====


% ===== СЕКЦИЯ 4: Статистика Ферми-Дирака =====
\section{Статистика Ферми-Дирака}

\subsection{Принцип Паули для фермионов}

Мы знаем замечательное математическое свойство: экспонента от суммы равна произведению экспонент. Применим это к нашему выражению:
\begin{equation}
    \mathcal{Z} = \sum_{\{n_{p\sigma}\}} \prod_{\vec{p}, \sigma} \exp\left( \frac{(\mu - \epsilon_{p\sigma}) n_{p\sigma}}{T} \right)
\end{equation}
Здесь идет произведение по всем возможным значениям импульса $p$ и проекции спина $\sigma$. 

А дальше мы можем поменять местами суммирование и произведение. Поскольку суммирование идет по всем возможным наборам $\{n_{p\sigma}\}$, независимое суммирование по каждому $n_{p\sigma}$ можно внести внутрь произведения. В каждом множителе теперь идет своя сумма по всем возможным значениям числа заполнения для данного конкретного состояния:
\begin{equation}
    \mathcal{Z} = \prod_{\vec{p}, \sigma} \sum_{n_{p\sigma}} \exp\left( \frac{(\mu - \epsilon_{p\sigma}) n_{p\sigma}}{T} \right)
\end{equation}

Какие значения принимает $n_{p\sigma}$? В силу принципа Паули для фермионов, на одном квантовом состоянии может находиться не более одной частицы, поэтому $n_{p\sigma}$ принимает только два значения: 0 и 1.

\subsection{Вычисление статсуммы для фермионов}

Поскольку $n_{p\sigma} \in \{0, 1\}$, сумма внутри произведения состоит всего из двух слагаемых. 
Если частиц нет ($n_{p\sigma}=0$), показатель экспоненты равен нулю, и экспонента дает единицу. Если есть один фермион ($n_{p\sigma}=1$), мы получаем саму экспоненту.

Подставляя эти два значения, мы находим Большую статистическую сумму для идеального ферми-газа:
\begin{equation}
    \mathcal{Z}_{FD} = \prod_{\vec{p}, \sigma} \left( 1 + \exp\left( \frac{\mu - \epsilon_{p\sigma}}{T} \right) \right)
\end{equation}
Вот мы и посчитали Большую статсумму!

\subsection{Средние числа заполнения: Распределение Ферми-Дирака}

Давайте сразу посмотрим на среднее число частиц в системе. В Большом каноническом ансамбле мы знаем удобную формулу для среднего числа частиц через производную логарифма статсуммы по химическому потенциалу:
\begin{equation}
    \langle N \rangle = T \frac{\partial \ln \mathcal{Z}}{\partial \mu}
\end{equation}

Для начала найдем логарифм нашей статсуммы. Логарифм произведения равен сумме логарифмов:
\begin{equation}
    \ln \mathcal{Z} = \sum_{\vec{p}, \sigma} \ln\left( 1 + \exp\left( \frac{\mu - \epsilon_{p\sigma}}{T} \right) \right)
\end{equation}

Теперь вычисляем производную:
\begin{equation}
    \langle N \rangle = T \frac{\partial}{\partial \mu} \sum_{\vec{p}, \sigma} \ln\left( 1 + \exp\left( \frac{\mu - \epsilon_{p\sigma}}{T} \right) \right)
\end{equation}
Дифференцируем логарифм. Производная от логарифма дает единицу, деленную на аргумент, а производная от внутренней экспоненты по $\mu$ «спускает» множитель $1/T$:
\begin{equation}
    \langle N \rangle = T \sum_{\vec{p}, \sigma} \frac{1}{1 + \exp\left( \frac{\mu - \epsilon_{p\sigma}}{T} \right)} \cdot \exp\left( \frac{\mu - \epsilon_{p\sigma}}{T} \right) \cdot \frac{1}{T}
\end{equation}

Температура $T$ перед суммой и в знаменателе благополучно сокращается. Чтобы привести выражение к стандартному виду, домножим числитель и знаменатель дроби на $\exp\left(-\frac{\mu - \epsilon_{p\sigma}}{T}\right)$ (или, что то же самое, просто перенесем экспоненту из числителя в знаменатель):
\begin{equation}
    \langle N \rangle = \sum_{\vec{p}, \sigma} \frac{1}{\exp\left( \frac{\epsilon_{p\sigma} - \mu}{T} \right) + 1}
\end{equation}

С другой стороны, мы с вами изначально писали, что полное среднее число частиц равно сумме средних чисел заполнения каждого состояния:
\begin{equation}
    \langle N \rangle = \sum_{\vec{p}, \sigma} \langle n_{p\sigma} \rangle
\end{equation}

Сравнивая эти два выражения (почленно для каждого состояния), мы получаем фундаментальный результат для среднего числа частиц на заданном энергетическом уровне:
\begin{equation}
    \langle n_{p\sigma} \rangle = \frac{1}{\exp\left( \frac{\epsilon_{p\sigma} - \mu}{T} \right) + 1}
\end{equation}
Это и есть знаменитая функция распределения Ферми-Дирака! Обычно мы будем обозначать её как $n_F(\epsilon_{p\sigma})$. То, что вы по физическому смыслу изначально закладывали для фермионов, вы сейчас строго формально получили из аппарата Большого канонического ансамбля.

% ===== КОНЕЦ СЕКЦИИ 4 =====






% ===== СЕКЦИЯ 5: Переход от суммирования к интегрированию =====
\section{Переход от суммирования к интегрированию}

Следующий пункт нашего аппарата — это вычисление сумм по одночастичным состояниям. Нам надо научиться переходить от дискретных сумм к непрерывным интегралам, так как в макроскопических системах уровни энергии расположены исключительно плотно.

\subsection{Интегрирование по импульсному пространству}

Предположим, у нас есть сумма по импульсам $\vec{p}$ от некоторой функции, зависящей от импульса:
\begin{equation}
    \sum_{\vec{p}} f(\vec{p})
\end{equation}

Как мы ее вычисляем? Вспомним, что на макроскопический ящик с газом (объемом $V = L^3$) наложены периодические граничные условия (условия Борна-Кармана). Это означает, что проекции импульса частицы $p_x, p_y, p_z$ принимают квазинепрерывные дискретные значения:
\begin{equation}
    p_{x,y,z} = \frac{2\pi\hbar}{L} n_{x,y,z}
\end{equation}
где $n_{x,y,z}$ — целые числа (положительные, отрицательные и ноль).

Какой объем в импульсном пространстве занимает одно такое квантовое состояние? Мы можем посчитать приращение импульса при изменении квантового числа на единицу ($\Delta n = 1$). Объем одного состояния $\Delta \Gamma_p$ равен:
\begin{equation}
    \Delta p_x \Delta p_y \Delta p_z = \left(\frac{2\pi\hbar}{L}\right)^3 \Delta n_x \Delta n_y \Delta n_z = \frac{(2\pi\hbar)^3}{V}
\end{equation}

Поскольку объем одного состояния равен $\frac{(2\pi\hbar)^3}{V}$, мы можем заменить суммирование по дискретным состояниям интегрированием по всему импульсному пространству, разделив элемент объема $d^3p$ на объем одного состояния:
\begin{equation}
    \sum_{\vec{p}} f(\vec{p}) = \int \frac{d^3p}{\frac{(2\pi\hbar)^3}{V}} f(\vec{p}) = \frac{V}{(2\pi\hbar)^3} \int f(\vec{p}) d^3p
\end{equation}
Это фундаментальное правило перехода от суммирования по импульсам к интегрированию в статистической физике.

\subsection{Плотность одночастичных состояний $\nu(\epsilon)$}

Очень часто нам нужно суммировать функции, которые зависят от импульса только через кинетическую энергию, то есть функции вида $f(\epsilon_{\vec{p}})$. В этом случае гораздо удобнее перейти к интегрированию по энергии. Для этого мы вводим плотность одночастичных состояний $\nu(\epsilon)$. 

Мы хотим переписать интеграл по импульсу в виде интеграла по энергии:
\begin{equation}
    \sum_{\vec{p}} f(\epsilon_{\vec{p}}) = \frac{V}{(2\pi\hbar)^3} \int f(\epsilon_{\vec{p}}) d^3p = \int_0^\infty \nu(\epsilon) f(\epsilon) d\epsilon
\end{equation}
Величина $\nu(\epsilon)$ по своему физическому смыслу показывает количество квантовых состояний системы, приходящихся на единичный интервал энергий вблизи значения $\epsilon$.

Давайте найдем явный вид этой функции. Для свободного газа (изотропный случай) элемент объема импульсного пространства в сферических координатах равен $d^3p = 4\pi p^2 dp$. Мы знаем связь энергии и импульса:
\begin{equation}
    \epsilon = \frac{p^2}{2m} \implies p = \sqrt{2m\epsilon}
\end{equation}
Берем дифференциал от энергии:
\begin{equation}
    d\epsilon = \frac{p}{m} dp \implies dp = \frac{m}{p} d\epsilon
\end{equation}
Подставляем всё в элемент объема $d^3p$:
\begin{equation}
    d^3p = 4\pi p^2 \left( \frac{m}{p} d\epsilon \right) = 4\pi m p d\epsilon = 4\pi m \sqrt{2m\epsilon} d\epsilon = 2\pi (2m)^{3/2} \sqrt{\epsilon} d\epsilon
\end{equation}

Отсюда мы считываем выражение для плотности состояний. Важно не забыть, что в реальных системах (например, для электронов) частицы обладают спином. Если функция $f(\epsilon)$ не зависит от проекции спина, то суммирование по спину $\sigma$ даст просто дополнительный множитель $g$ — фактор спинового вырождения (для электрона со спином $1/2$ этот фактор равен $g=2$).
Включая фактор $g$ в плотность состояний, получаем итоговую формулу:
\begin{equation}
    \nu(\epsilon) = g \frac{V}{(2\pi\hbar)^3} 2\pi (2m)^{3/2} \sqrt{\epsilon} = \frac{g V m^{3/2}}{\sqrt{2} \pi^2 \hbar^3} \sqrt{\epsilon}
\end{equation}

Как мы видим, для трехмерного идеального газа плотность состояний пропорциональна корню из энергии:
\begin{equation}
    \nu(\epsilon) = A \sqrt{\epsilon}
\end{equation}
где $A$ — константа, зависящая от массы частиц и объема системы. 

\begin{figure}[h!]
    \centering
    \includegraphics[width=0.3\textwidth]{"images/nu(e).png"}
    \caption{Плотность одночастичных состояний идеального трехмерного газа.}
\end{figure}

Обязательно запомните этот факт: $\nu(\epsilon) \propto \sqrt{\epsilon}$. Для других размерностей пространства или для другого спектра (например, фотонов или фононов) зависимость будет другой, но для нашего классического спектра $\epsilon = p^2/2m$ в трехмерном случае это всегда корень.

\subsection{Интегральные выражения для $\langle N \rangle$ и $\langle E \rangle$}

Следовательно, теперь все наши термодинамические средние величины можно записать в виде интегралов. Мы рассматриваем основной случай — когда ферми-газ не находится во внешнем магнитном поле, то есть энергия одночастичного уровня от проекции спина $\sigma$ не зависит: $\epsilon_{p\sigma} = \epsilon_p$. Тогда суммирование по $\sigma$ независимо и просто даёт кратность спинового вырождения $g$ (для электронов $g = 2$). Этот множитель уже включён в нашу плотность состояний $\nu(\epsilon)$, поэтому любая сумма по $\vec{p}$ и $\sigma$ от функции, зависящей только от энергии, превращается в интеграл по энергии:
\begin{equation}
    \sum_{\vec{p}, \sigma} f(\epsilon_p) = \int_0^\infty \nu(\epsilon)\, f(\epsilon)\, d\epsilon
\end{equation}

Применяем это правило к нашим термодинамическим средним. Для среднего числа частиц:
\begin{equation}
    \langle N \rangle = \sum_{\vec{p}, \sigma} \langle n_{p\sigma} \rangle = \int_0^\infty \nu(\epsilon)\, n_F(\epsilon)\, d\epsilon
\end{equation}
Для средней энергии:
\begin{equation}
    \langle E \rangle = \sum_{\vec{p}, \sigma} \epsilon_p\, \langle n_{p\sigma} \rangle = \int_0^\infty \epsilon\, \nu(\epsilon)\, n_F(\epsilon)\, d\epsilon
\end{equation}
где $n_F(\epsilon) = \dfrac{1}{e^{(\epsilon - \mu)/T} + 1}$ — функция распределения Ферми-Дирака.

Вот это и есть основные уравнения задачи. Именно такие интегралы нам предстоит научиться вычислять. Но прежде чем переходить к их вычислению (а аналитически они точно не берутся в общем случае и потребуют перехода к какому-то конкретному пределу), давайте разберёмся, когда реализуются классический и квантовый режимы.


% ===== СЕКЦИЯ 6: Классический и квантовый пределы ферми-газа =====
\section{Классический и квантовый пределы ферми-газа}

\subsection{Классический предел: переход к распределению Больцмана}

Первый очевидный случай, который мы должны хорошо понять, — это переход к классическому газу. Когда ферми-газ можно рассматривать как классический газ Больцмана?

Классический предел реализуется, когда состояний намного больше, чем частиц. Это значит, что в среднем на каждом состоянии находится намного меньше одной частицы:
\begin{equation}
    \langle n_{p\sigma} \rangle \ll 1
\end{equation}
В этом случае функция Ферми-Дирака мала, что соответствует условию, что экспонента в знаменателе намного больше единицы:
\begin{equation}
    e^{(\epsilon - \mu)/T} \gg 1
\end{equation}
Это выполняется, когда $e^{-\mu/T} \gg 1$, то есть химический потенциал $\mu$ достаточно отрицателен.

В этом пределе $+1$ в знаменателе Ферми-Дирака можно отбросить, и функция распределения переходит в распределение Больцмана:
\begin{equation}
    n_F(\epsilon) = \frac{1}{e^{(\epsilon - \mu)/T} + 1} \approx e^{(\mu - \epsilon)/T}
\end{equation}

\subsection{Вычисление $\langle N \rangle$ в классическом пределе}

Подставим приближение Больцмана в интегральное выражение для среднего числа частиц:
\begin{equation}
    \langle N \rangle = \int_0^\infty \nu(\epsilon)\, e^{(\mu - \epsilon)/T}\, d\epsilon = A\, e^{\mu/T} \int_0^\infty \sqrt{\epsilon}\, e^{-\epsilon/T}\, d\epsilon
\end{equation}
Делаем замену переменной $x = \epsilon / T$:
\begin{equation}
    \langle N \rangle = A\, e^{\mu/T}\, T^{3/2} \int_0^\infty \sqrt{x}\, e^{-x}\, dx
\end{equation}
Стоящий здесь интеграл — это по определению гамма-функция $\Gamma(3/2)$. Значение $\Gamma(3/2) = \frac{\sqrt{\pi}}{2}$. Следовательно:
\begin{equation}
    \langle N \rangle = A\, T^{3/2}\, \Gamma\!\left(\tfrac{3}{2}\right) e^{\mu/T}
\end{equation}

Подставим явный вид константы $A$ и приведём выражение к удобному виду. После подстановки обнаруживается, что в результате образуется тепловая длина волны де Бройля в кубе — квантовый объём $V_Q$:
\begin{equation}
    V_Q = \left( \frac{2\pi\hbar^2}{mT} \right)^{3/2}
\end{equation}
Тогда:
\begin{equation}
    \langle N \rangle = \frac{g V}{V_Q}\, e^{\mu/T}
\end{equation}

\subsection{Химический потенциал классического газа}

Из полученного выражения легко найти химический потенциал:
\begin{equation}
    e^{\mu/T} = \frac{\langle N \rangle\, V_Q}{g V} = \frac{n\, V_Q}{g}
\end{equation}
\begin{equation}
    \mu = T \ln\!\left( \frac{n\, V_Q}{g} \right)
\end{equation}
где $n = \langle N \rangle / V$ — концентрация частиц.

В классическом режиме число частиц в квантовом объёме намного меньше единицы: $n V_Q \ll 1$. Это и есть критерий классического газа. Из него сразу следует, что $\mu < 0$, как и полагается для классического газа. Именно такую формулу для химического потенциала мы уже получали на второй лекции — смотрите, она сама собой воспроизводится из аппарата Большого канонического ансамбля.

\subsection{Температура вырождения}

Итак, данная формула для химического потенциала справедлива при температурах намного больше температуры вырождения:
\begin{equation}
    T \gg T_{\text{выр}}
\end{equation}
Температура вырождения — это та температура, при которой в квантовом объёме $V_Q(T)$ оказывается порядка одной частицы:
\begin{equation}
    n\, V_Q(T_{\text{выр}}) \sim 1
\end{equation}
Из этого условия, подставив явный вид $V_Q$, мы оцениваем температуру вырождения:
\begin{equation}
    T_{\text{выр}} \sim \frac{\hbar^2}{m}\, n^{2/3}
\end{equation}

Для нормального металла — например, меди — концентрация электронов проводимости $n \approx 5 \cdot 10^{22}$~см$^{-3}$, масса — масса электрона $m = m_e$. Подставляем и получаем:
\begin{equation}
    T_{\text{выр}} \approx 5 \cdot 10^4 \text{ К}
\end{equation}
Это колоссальная величина. При такой температуре от медного куска ничего не осталось бы — она давно испарилась. При комнатной температуре ($T \sim 300$~К) условие $T \ll T_{\text{выр}}$ выполнено с огромным запасом, поэтому нас в первую очередь интересует квантовый режим.

\subsection{Ультраквантовый предел ($T = 0$): энергия Ферми}

Давайте рассмотрим ультраквантовый случай $T \to 0$. При нулевой температуре функция Ферми-Дирака вырождается в ступенчатую функцию Хевисайда: она равна единице при $\epsilon \leq \mu(0) \equiv \epsilon_F$ и нулю при $\epsilon > \epsilon_F$:
\begin{equation}
    n_F(\epsilon) \xrightarrow{T \to 0} \begin{cases} 1, & \epsilon \leq \epsilon_F \\ 0, & \epsilon > \epsilon_F \end{cases}
\end{equation}

Здесь $\epsilon_F$ — \textbf{энергия Ферми}: значение химического потенциала $\mu$ при $T = 0$. Почему именно так? Потому что функция Ферми-Дирака при $T = 0$ равна единице для всех $\epsilon < \mu$ и нулю для $\epsilon > \mu$ — то есть ведёт себя ровно так, как мы определяли уровень Ферми в курсе квантовой механики. Важно учтите: $\epsilon_F$ — не абсолютная константа, она зависит от плотности частиц.

\subsection{Число частиц и импульс Ферми при $T = 0$}

Вычислим число частиц в системе при $T = 0$ через интеграл по энергии. Ступенчатая функция Хевисайда обрезает интеграл на уровне $\epsilon_F$:
\begin{equation}
    N = \int_0^{\epsilon_F} \nu(\epsilon)\, d\epsilon = A \int_0^{\epsilon_F} \sqrt{\epsilon}\, d\epsilon = A \cdot \frac{2}{3}\, \epsilon_F^{3/2}
\end{equation}
Замечаем, что $A \sqrt{\epsilon_F} = \nu(\epsilon_F)$ — плотность состояний на поверхности Ферми. Поэтому:
\begin{equation}
    N = \frac{2}{3}\, \epsilon_F\, \nu(\epsilon_F)
\end{equation}
Запомните эту полезную формулу: число частиц выражается через энергию Ферми и плотность состояний на уровне Ферми.

Теперь вычислим $N$ альтернативным способом — непосредственно в пространстве импульсов. При $T = 0$ все состояния с $|\vec{p}| \leq p_F$ (внутри «ферми-сферы») заполнены единицей, а с $|\vec{p}| > p_F$ — пусты. Поэтому:
\begin{equation}
    N = \frac{g V}{(2\pi\hbar)^3} \int_{|\vec{p}| \leq p_F} d^3p = \frac{g V}{(2\pi\hbar)^3} \cdot \frac{4\pi}{3}\, p_F^3
\end{equation}
Для электронов с $g = 2$:
\begin{equation}
    N = \frac{V}{3\pi^2 \hbar^3}\, p_F^3
\end{equation}
Отсюда выражаем импульс Ферми через концентрацию $n = N/V$:
\begin{equation}
    p_F = \hbar\!\left( 3\pi^2 n \right)^{1/3}
\end{equation}
Через волновой вектор Ферми $k_F = p_F / \hbar$:
\begin{equation}
    k_F = \left( 3\pi^2 n \right)^{1/3}
\end{equation}

\subsection{Энергия Ферми и её связь с температурой вырождения}

Из дисперсионного соотношения $\epsilon_F = p_F^2 / (2m)$ находим энергию Ферми:
\begin{equation}
    \epsilon_F = \frac{\hbar^2 k_F^2}{2m} = \frac{\hbar^2}{2m}\!\left( 3\pi^2 n \right)^{2/3}
\end{equation}
По порядку величины:
\begin{equation}
    \epsilon_F \sim \frac{\hbar^2}{m}\, n^{2/3}
\end{equation}
Это в точности то же самое выражение, что и для температуры вырождения. Именно поэтому $T_{\text{выр}} \sim \epsilon_F$, и эти понятия часто используются как синонимы. Для нормальных металлов (медь, алюминий и т.\,д.) энергия Ферми лежит в диапазоне:
\begin{equation}
    \epsilon_F \approx 1 \div 7 \text{ эВ}
\end{equation}
что соответствует температурам $T_{\text{выр}} \approx (1\div 7) \cdot 10^4$~К.

% ===== КОНЕЦ СЕКЦИИ 6 =====


% ===== СЕКЦИЯ 7: Критерий применимости модели идеального ферми-газа =====
\section{Критерий применимости модели идеального ферми-газа}

\subsection{Когда можно пренебречь кулоновским взаимодействием?}

Следующий пункт: когда можно пользоваться приближением невзаимодействующего идеального ферми-газа? Мы говорили, что фермионы не взаимодействуют и заполняют одночастичные уровни энергии. Но электроны заряжены и взаимодействуют по закону Кулона — когда это взаимодействие можно пренебречь?

Критерий применимости идеального газа: средняя кинетическая энергия на один электрон должна быть намного больше кулоновской энергии взаимодействия электрона со своим ближайшим соседом:
\begin{equation}
    \bar{E}_{\text{кин}} \gg U_{\text{Кул}} \sim \frac{e^2}{\bar{r}}
\end{equation}
где $\bar{r}$ — среднее расстояние между электронами.

Оцениваем каждую из величин. Среднее расстояние между электронами связано с концентрацией очевидным образом:
\begin{equation}
    \bar{r} \sim \left(\frac{V}{N}\right)^{1/3} = \frac{1}{n^{1/3}}
\end{equation}

Средняя кинетическая энергия электронов в вырожденном ферми-газе по порядку величины равна энергии Ферми:
\begin{equation}
    \bar{E}_{\text{кин}} \sim \epsilon_F \sim \frac{\hbar^2}{m}\, n^{2/3}
\end{equation}

Подставляем в условие применимости:
\begin{equation}
    \frac{\hbar^2}{m}\, n^{2/3} \gg e^2 n^{1/3}
\end{equation}
Сокращаем на $n^{1/3}$ и переносим $e^2$ в левую часть:
\begin{equation}
    \frac{\hbar^2}{m e^2}\, n^{1/3} \gg 1
\end{equation}
Дробь $\hbar^2 / (m e^2)$ — это в точности боровский радиус $a_B$. Поэтому:
\begin{equation}
    a_B\, n^{1/3} \gg 1 \qquad \text{или, что то же самое,} \qquad a_B \gg \bar{r}
\end{equation}
Через импульс Ферми это эквивалентное условие можно записать ещё так:
\begin{equation}
    a_B\, k_F \gg 1
\end{equation}

\subsection{Парадоксальный результат: плотный газ — идеальный газ}

Что означает условие $a_B \gg \bar{r}$? Среднее расстояние между электронами должно быть намного \emph{меньше} боровского радиуса. Другими словами, газ должен быть очень \emph{плотным}! Это совершенно парадоксальный и контринтуитивный результат.

Физически это объясняется следующим образом. В вырожденном ферми-газе кинетическая энергия растёт с плотностью как $\epsilon_F \sim n^{2/3}$, тогда как кулоновская энергия взаимодействия растёт медленнее: $e^2/\bar{r} \sim e^2 n^{1/3}$. Поскольку кинетическая энергия при увеличении плотности растёт быстрее потенциальной, в пределе высокой плотности кинетика подавляет взаимодействие — и газ становится идеальным. Для классического газа всё наоборот.

У реальных металлов среднее расстояние $\bar{r}$ имеет порядок боровского радиуса $a_B$ или чуть больше. То есть условие идеальности для металлов выполняется на пределе, но не с большим запасом. Тем не менее, даже в этом случае модель идеального ферми-газа служит превосходным первым приближением и позволяет качественно и во многом количественно описывать свойства металлов.

Вот мы сейчас только разобрали «вершки» — саму основу аппарата. Нам необходимо ещё научиться брать интегралы, содержащие ферми-функцию распределения при конечных низких температурах. Пока мы лишь на подступах. В следующий раз мы как раз с этого и начнём: возьмём эти интегралы и получим явные температурные зависимости химического потенциала, энергии и теплоёмкости ферми-газа при низких температурах. На сегодня всё.

% ===== КОНЕЦ СЕКЦИИ 7 =====






\end{document}

